\section{Design Reviews Before Coding (implementation\_plan.md)}

This looks at what the plan document itself is actually required to contain, and in particular at the Design Review portion, which is the part of the document the human review gate is centered on.

An \textbf{\texttt{implementation\_plan.md}} follows a fixed eight-part structure: a feature summary, the mandatory Design Review, a list of files to create, a list of files to modify, any API contract changes, the corresponding OpenAPI specification additions, an explicit account of any deviation from the charter's rules, and a verification plan describing how to confirm the feature actually works. Of these, only the Design Review is treated as non-negotiable --- a plan missing it, or missing any of its five required sub-sections, is considered invalid and cannot be approved.

Those five sub-sections are: an architecture rationale, explaining why a particular file or component structure was chosen over the alternatives and which existing pattern it follows; a statement of state ownership, identifying, for every new piece of state the feature introduces, which component or hook owns it and why; a statement of service ownership, identifying which service file is responsible for which piece of logic and justifying whether an existing service is being extended or a new one created; a testing strategy, specifying what will be tested, at what level, and which failure scenarios are explicitly covered; and a set of scalability concerns, naming any decision that could become a bottleneck or a maintenance burden as the feature or its data grows, along with a proposed mitigation. The purpose of requiring all five before any code is written is to force the reasoning behind a design to be made explicit up front, rather than letting a plan consist of little more than a list of files to touch.

A concrete instance of this is the plan produced for adding a super-admin role to the application. Before any code existed, the plan flagged, under a dedicated review section, that the feature would require two new columns on the users table, two new database tables for audit logging and user notifications, a routing guard restricting the new admin area to admin accounts, and --- notably --- a privacy concern: since a super admin would be able to browse other users' saved requests and execution histories, the plan proposed a visible warning banner in that view and required every instance of an admin viewing another user's data to be recorded immutably in the audit log. In its Design Review, the plan gave each of the five required sub-sections a concrete answer for this feature: the architecture rationale justified isolating the new administrative logic under its own backend blueprint and service module rather than extending an existing one; state ownership was laid out in a table assigning each new piece of state --- the admin flag itself, the active tab in the admin panel, the users and workspaces lists, in-app notifications --- to the specific part of the application responsible for it; and service ownership separated the backend service handling admin actions from the frontend service wrapping the corresponding API calls.

What this example shows is what a mandatory design review actually buys in practice: a privacy concern with legal and ethical weight was identified and a mitigation proposed while the feature still existed only as a document, reviewable and rejectable by a human before a single line of code had been written --- rather than being discovered only after the fact, in the way several of Part 1's problems were.

